\section{Empirical Strategy}

The empirical objective is to identify the dynamic causal pass-through from tradable
inflation, $\pi^T_{it}$, to non-tradable inflation, $\pi^N_{it}$, in a monthly panel of
19 euro area countries.
Because tradable and non-tradable inflation comove with common demand conditions and
aggregate cost pressures, ordinary least squares conflates pass-through with spurious
comovement; identification requires an external source of variation in $\pi^T_{it}$
that is orthogonal to domestic non-tradable conditions.
The instrument is a Bartik shift-share of the form $Z_{it} = w_{it} \cdot Z_t$, where
$w_{it}$ is the tradables expenditure share of country $i$'s CPI basket in month $t$
--- computed as the sum of ECOICOP expenditure weights for food, clothing, furnishings,
transport, and alcohol --- and $Z_t$ is the market-based oil price expectation shock of
\citet{baumeister2019}.
Heterogeneity in $w_{it}$ across countries and over time ensures that the same common
oil shock translates into proportionally larger cost impulses for countries with larger
tradables baskets, generating cross-sectional identifying variation that is unrelated to
idiosyncratic domestic conditions.
The exclusion restriction requires that $Z_{it}$ affects non-tradable prices only through
its effect on tradable inflation; any direct transmission through energy costs in service
production is absorbed by twelve lags of both inflation series and the unemployment rate,
and country fixed effects absorb permanent cross-country differences in non-tradable price
setting.
Time fixed effects are excluded because they would absorb $Z_t$, eliminating the
instrument.

The dynamic pass-through is estimated horizon by horizon following the local projection
method of \citet{jorda2005}.
The first stage establishes instrument relevance by regressing tradable inflation on $Z_{it}$
and controls:
\begin{equation}
    \pi^T_{it}
    = \alpha_i
    + \lambda\, Z_{it}
    + \sum_{\ell=1}^{L} \boldsymbol{\gamma}_\ell'\, \mathbf{x}_{i,t-\ell}
    + v_{it},
    \label{eq:firststage}
\end{equation}
where $\alpha_i$ are country fixed effects, $Z_{it} = w_{it}\cdot Z_t$ is the
Bartik instrument, and
$\mathbf{x}_{i,t-\ell} = (\pi^N_{i,t-\ell},\,\pi^T_{i,t-\ell},\,u_{i,t-\ell})'$
collects lags of non-tradable inflation, tradable inflation, and the unemployment rate.
The second stage projects non-tradable inflation $h$ months ahead onto instrumented
tradable inflation:
\begin{equation}
    \pi^N_{i,t+h}
    = \alpha_i
    + \theta_h\,\widehat{\pi^T}_{it}
    + \sum_{\ell=1}^{L} \boldsymbol{\gamma}_\ell'\, \mathbf{x}_{i,t-\ell}
    + \varepsilon_{i,t+h},
    \quad h = 0, 1, \ldots, H,
    \label{eq:lpiv}
\end{equation}
where $\widehat{\pi^T}_{it}$ denotes tradable inflation instrumented by $Z_{it}$.
The sequence $\{\theta_h\}_{h=0}^{H}$ traces the dynamic second-round pass-through from
tradable to non-tradable inflation, with cumulative pass-through at horizon $H$ given by
$\sum_{h=0}^{H}\theta_h$.
The main specification uses $H = 13$ monthly horizons and $L = 12$ lags; standard errors
are robust to heteroskedasticity throughout.

To assess whether pass-through varies across economic states, equation~\eqref{eq:lpiv}
is estimated separately for subsamples defined by regime indicators, an approach that
avoids the need for an additional excluded instrument to identify the triple interaction
$\widehat{\pi^T}_{it} \cdot D_{it}$.
Three regime dimensions are considered.
First, the sample is split according to whether the Global Supply Chain Pressure Index
(GSCPI) of \citet{benigno2022} is above or below its standardised long-run mean of zero,
separating episodes of supply-chain stress from normal logistical conditions.
Second, the sample is split by the sign of the log monthly change in the nominal effective
exchange rate, $\Delta\ln\mathrm{NEER}_{it}$, testing for asymmetric pass-through between
appreciation and depreciation episodes.
Third, the inflation environment is defined by whether lagged headline HICP inflation
exceeds the ECB's 2 percent target; using the lagged value ensures the regime indicator
is predetermined with respect to the outcome.
Each subsample retains the same instrument, controls, and fixed effects as the baseline
specification.
The full sample spans January 1998 to December 2025, with pre- and post-2020 sub-periods
reported to assess parameter stability around the pandemic energy-price episode.
